\section{GIỚI THIỆU}
\subsection{Tổng quan}
Đề tài "Thiết kế máy chơi game cầm tay BKONSOLE" đặt ra các nhiệm vụ cụ thể về cả phần cứng và phần mềm nhằm xây dựng một hệ thống nhúng hoàn chỉnh, hoạt động ổn định và đáp ứng nhu cầu giải trí cơ bản.
\subsection{Nhiệm vụ đề tài}
\textit{Nhiệm vụ cốt lõi của phần cứng là tạo ra một mạch điện tử tích hợp, nhỏ gọn và đảm bảo các khối chức năng hoạt động đồng bộ với vi điều khiển trung tâm. Cụ thể:}
\begin{itemize}
    \item Khối Xử lý Trung tâm: Thiết kế mạch nguyên lý và PCB cho vi điều khiển ATmega328P (hoạt động ở 16MHz), đảm bảo kết nối chính xác các chân chức năng (GPIO, SPI, UART) và khả năng nạp chương trình.
    \item Thiết kế khối hiển thị: Kết nối và điều khiển Module LED Matrix 8x8 thông qua IC ghi dịch MAX7219 (giao tiếp SPI) để hiển thị giao diện chính của trò chơi (rắn săn mồi, hứng gạch, v.v.). Yêu cầu hiển thị rõ ràng, không bị nhấp nháy hay bóng mờ. Ngoài ra tích hợp màn hình LCD 1602 để hiển thị các thông tin phụ trợ như: tên trò chơi, điểm số (Score), trạng thái (Game Over/Start), hoặc menu cài đặt.
    \item Thiết kế khối nguồn năng lượng: Xây dựng hệ thống cấp nguồn di động sử dụng pin Li-Ion/Li-Po 3.7V. Hệ thống phải bao gồm:
        \begin{itemize}
            \item Mạch sạc pin tích hợp bảo vệ sử dụng IC TP4056 để đảm bảo an toàn khi sạc qua cổng USB.
            \item Mạch nâng áp (Boost Converter) sử dụng IC MT3608 để ổn định điện áp đầu ra ở mức 5V cấp cho vi điều khiển và các module hiển thị, đảm bảo hệ thống không bị sụt áp khi tải tăng cao.
        \end{itemize}
    \item Thiết kế khối điều khiển: Bố trí các nút nhấn (Push Buttons) hợp lý để người dùng thao tác (Di chuyển Lên, Xuống, Trái, Phải và nút Chọn/Start).
\end{itemize}

\textit{Phần mềm đóng vai trò điều phối hoạt động của toàn bộ hệ thống. Nhiệm vụ lập trình bao gồm:}
\begin{itemize}
    \item Xây dựng Thư viện Giao tiếp (Driver): Lập trình các hàm điều khiển để ATmega328P giao tiếp với ngoại vi, đặc biệt là sử dụng giao thức SPI cho LED Matrix (MAX7219) và điều khiển LCD 1602.
    \item Xây dựng Logic Trò chơi (Game Logic): Lập trình thuật toán cho các game cổ điển (Rắn săn mồi, Pong,...) trên ma trận 8x8. Xử lý các tác vụ thời gian thực như quét nút nhấn, tính toán va chạm, và cập nhật điểm số.
    \item Quản lý Hệ thống: Sử dụng Timer và Interrupt của vi điều khiển để đảm bảo game chạy mượt mà, tốc độ phản hồi nhanh và kiểm soát thời gian làm tươi màn hình.
\end{itemize}